\section{Discussion}

\towrite[Aim: open the chapter in one short paragraph. Cover: how it interprets the results against the research question.]

\subsection{Introduction}

\towrite[Aim: outline the chapter. Cover: restate the research question and the order of the sections that follow.]

\subsection{Interpretation of Experimental Results}

\towrite[Aim: say what the results mean. Cover: the main findings and their likely explanations; unexpected results; whether each hypothesis held and why. Distinguish what the data show from what you infer.]

\subsection{Evaluation against Research Objectives}

\towrite[Aim: assess the work against its objectives. Cover: for each objective in Section 1.7, the supporting evidence; a table may help. Answer the research question directly.]

\subsection{Reliability of BCI-Based Car Control}

\towrite[Aim: judge whether commands are dependable enough. Cover: what success, missed command and false activation rates mean for real use; which factors affected reliability, such as user, command or training.]

\subsection{Latency and Suitability for Real-Time Control}

\towrite[Aim: judge whether the system is fast enough. Cover: what the delays mean for driving a car; which stage contributes most delay; which parts could improve and which depend on Emotiv's detection.]

\subsection{Trade-offs between Speed and Reliability}

\towrite[Aim: discuss the central design tension. Cover: how stricter thresholds and smoothing reduce errors but add delay; what your data shows about this; the balance you recommend, and for what use.]

\subsection{Implications for Consumer-Grade EEG Applications}

\towrite[Aim: look beyond this car. Cover: what the results suggest for assistive technology, robotics and other uses of consumer EEG; where such devices are and are not appropriate, including safety-critical limits.]

\subsection{Comparison with Previous Studies}

\towrite[Aim: place results in the literature. Cover: comparison with the studies in Section 2.9, by device, paradigm and metric; agreements and differences; why direct comparison may be limited.]

\subsection{Technical and Methodological Limitations}

\towrite[Aim: assess the study honestly. Cover: sample size and participant mix; controlled environment; black-box detection; hardware and timing measurement limits; how each may have affected the results.]

\subsection{Potential Improvements and Future Work}

\towrite[Aim: suggest practical improvements to this system. Cover: changes to the framework, control logic, safety and protocol, each linked to a finding or limitation. Broader research directions belong in Section 7.6.]

\subsection{Chapter Summary}

\towrite[Aim: close the chapter. Cover: the main interpretation and its significance, leading into the conclusion.]
